Beekeeping in the semiarid region of Ceará, although advantageous, faces competitiveness barriers that associativism, grounded in the Solidarity Economy, helps to overcome. However, delegating duties to a board creates information asymmetry among members, which may undermine trust and group cohesion. This work aims to propose a web application to automate the internal processes of the Associação de Apicultores de Poranga, strengthening management transparency. Given the members' low familiarity with digital tools, requirements will be elicited through Exploratory Prototyping, conducted in iterative cycles with stakeholders representative of the main user profiles, with the prototype built following the Mobile First approach. Development will follow the Test-Driven Development (TDD) cycle, with automated tests (unit, integration, and regression) ensuring system reliability throughout the codebase's evolution, complemented by a usability evaluation conducted with members not involved in the elicitation stage. By the end of the process, the application is expected to be able to reliably automate the association's internal processes, making information available transparently and accessibly to all members, in order to reduce information asymmetry, strengthen mutual trust, and sustain the organization's collective self-management.
 
\textbf{Keywords:} Beekeeping. Solidarity Economy. Associativism. Software Engineering. Test-Driven Development.